\chapter{Introductory Double Descent}
\begin{description}
  \item[Entry] The chapter is of the entry for \textbf{Double Descent in Machine Learning - An overview} (Bui Gia Khanh, planned March 2026). 
\end{description}

The analysis of learning action, machine learning and related practices in the field theoretically has been made by utilizing Computational Learning Theory (CoLT) \cite{10.1145/1968.1972} and Statistical Learning Theory (SLT) \cite{Vapnik1999-VAPTNO}. These two theories, while overlapped, provides a general framework in analysing and justifying learning actions and learning model constructions, aiding in the formation of modern practices. One of the famous insight using such framework is the \textit{bias-variance tradeoff} \cite{6797087}, which states that model complexity and generality inversely affect each other, thus guarantee the need for a safe bracket between them. However, recent literatures, \cite{belkin_reconciling_2019} has indicated the fallout of such dilemma by the new phenomenon called \textit{double descent}, where there exists an interpolation threshold that renders the current statistical justification for bias-variance inconsequential to a given region. Various `anomaly' has also been detected similarly, with varying degree of sophistication and potential.In this paper, we would like to review through progress made in this particular field of research, and the phenomena potential explanations up until now. 

\section{Introduction}

The theory of machine learning and its modern practice has been developed and researched, of a substantial portion by empirical and heuristic approach, either by advancing new practices, architectures or by try-and-test modification. From its early onset of a regression estimator $\theta(x,y)$ on the linear regression problem, machine learning has developed substantially. Certain model concept with great successes includes regression-classification model, Bayesian modelling, generative learning model, support vector machine, Gaussian processes, and more. Of all such, the more formal and complex model architecture created, is the concept of a \textit{neural network}. With increasingly sophisticated architecture, heuristic approach become popular, the fast-paced advancement of the field comes with new method, new results, new observations, and its far-reaching application which led to even bigger and larger scale deployment, there has been questions about the formation and status of a theoretical ground, a rigorous matter on the side of \textbf{theoretical machine learning}.

While rigorous and well-formulated in a sense, classical and theoretical machine learning was dwarfed by the modern advancement of machine learning as a whole, leading to several anecdotal problems regarding the interpretation of phenomena, the re-evaluation of the theory to fit the more updated analysis, and explanation to more sophisticatedly designed system. This and many more, plus as present, many of such advancements and improvements are heuristic, and the general theory and conceptual understanding remain limited, led to the choice to often opt for analogies and empirical workaround. Ultimately, this resulted in multiple `failures' in explaining, interpreting, and predicting behaviours of new phenomena appeared in the modern landscape of machine learning. Thereby, we suggest some particular insights and analysis, and a fix within interpretation of the phenomena. 

Our main focus is to review on the topic of the umbrella term classical learning theory, and the recent phenomena observed named `double descent'. Statistical learning theory (SLT) and computational learning theory (CLT) \cite{Vapnik1999-VAPTNO,10.5555/2371238,10.5555/2621980,STL_Hajek_Maxim_2021,bousquet2020theoryuniversallearning} has been prominent in constructing a well-rounded formal theory surrounding learning problems, models, and machine learners analysis. Valuable insights have been dissected from treatment of statistical theory and mathematical modelling on models, including the \textit{bias-variance tradeoff} \cite{6797087,Domingos2000AUB}, which serves as a bound for efficient learning and model configuration (or complexity). However, recently there has been observations of \textit{double descent} \cite{belkin_reconciling_2019,schaeffer_double_2023,nakkiran_deep_2019,lafon_understanding_2024} which refute the famous tradeoff assumption, and hence brings question to the establishment of the theory, as well as several assumptions and insight in the framework. Further events and phenomena observed also includes grokking and triple, to $n$-descent \cite{davies_unifying_2023,d_ascoli_triple_2020}. Furthermore, many problems of defining and formalizing notions used in designing and implementing machine learning models are inconclusive, such as, for example, \textit{model complexity} and others. 

On itself, the phenomena \textit{double descent} has been investigated somewhat comprehensively, firstly introduced by \cite{belkin_reconciling_2019}, which gives it distinctive saddle escape illustration. Further analysis was made by several literatures, particularly attributed the existence of double descent to the concept of model complexity and inductive bias, and potential explanations of such. \cite{nakkiran_deep_2019} expanded the phenomena into deep neural network models. Their conclusion is reached by considering the perturbation of a learning procedure $\mathcal{T}$ on the effective model complexity $\mathrm{EMC}_{\mathcal{D},\epsilon}(\mathcal{T})$ defined in the paper, separating the eventual phenomena into regions of observations, thereby in one way or another, predicting the tendency of double descent. This is also the first, arguably, "concise" definition of double descent, even though its nature is an empirical definition, including the notion of model complexity. Recently, there is also the Neural Tangent Kernel (\cite{Jacot:2018:NTK}) which can be used to explain double descent, however further researches are still required. Preliminary works are done by \cite{lafon_understanding_2024}, \cite{schaeffer_double_2023}, and \cite{liu2023understandingroleoptimizationdouble} on the role of optimization in double descent. \cite{davies_unifying_2023} attempted to unifying grokking with double descent, theorized a possibility of similarity during the generalization phase transition of the inference period, and \cite{olmin2024understandingepochwisedoubledescent} attempted to explain epoch-wise double descent within the model of two-layer linear network. However, it is mentioned that it can also be expanded into deep nonlinear networks. 

\section{Problem statements}

Before continuing further, we must review the problem statement that led to the problem of double descent. We will have a look at the theory, the original conflict, and the illustrative phenomenon in action. 

\subsection{Bias-variance tradeoff}

In classical sources, \cite{gareth_james_introduction_2013,goodfellow2016deep,STL_Hajek_Maxim_2021,10.5555/2371238,10.5555/2621980}, bias-variance tradeoff is a classical principle used in the categorization of technique called \textit{model selection}. Specifically, this refers to the progress of choosing the best predictor $h^{*}$ that minimize the generalization error that is the main goal of a particular machine learning model, during the process of training and inference. Originally, the theory that bias and variance often come in to conjunction is \textbf{Estimation theory}, of which there exists an estimator $\hat{\theta}$ that use a set of observations, to estimate the concept, or the underlying mechanics of certain system that outputted the observation set, by the \textbf{probabilistic perspective} - that is, it is governed by appearance by an independently and identically distributed process of parameterized probability $\theta\in \Theta$. Here, parameterized means being expressed by a set, often finite, of parameters, by \cite{LehmannCasella_theory_1998,liam_statistics_2005}. It is only when \cite{6797087} introduced it to machine learning that we have an equivalent structure that is similar to bias-variance in machine learning. 

\subsubsection{Precursor (Geman et al., 1992)}

The original definition of bias-variance tradeoff by \cite{6797087} is first constructed using the means-square error, which is regarded as a normal measure in the real encoding space. Their approach is to justify bias-variance via decomposition of the loss function $\ell$, for such to find an alternative reasonable form of such loss landscape. Suppose of a regression problem to construct a hypothesis function $f(x)$ from $(x_{1},y_{1},\dots,x_{N},y_{N})$ for the purpose of generalization - that is, predicting unseen variational values for different pair $(x_{j}, \mathord{?})$ such that $\mathord{?}=y_{j}+\epsilon$ for a conceivable implicit error. To be explicit about the relation of this problem, or $f$ on the given data $\mathcal{D}=\{(x_i, y_i)\mid i \leq N\}$, denote $f(x;\mathcal{D})$ instead of $f$, the natural mean-square measure as a predictor is: 
\begin{equation}
    \mathcal{M}(f,y) = \mathbb{E} \left[((y-f(x;\mathcal{D})))^{2}\mid x, \mathcal{D}\right] 
\end{equation} for $\mathbb{E}[\cdot]$ the expectation wrt to a distribution $P$. Decomposing the right-hand side, we have: 
\begin{equation}
    \mathcal{M}(f,y) = \mathbb{E} \left[((y-f(x;\mathcal{D})))^{2}\mid x, \mathcal{D}\right] = \mathbb{E}\left[(y-\mathbb{E}[y\mid x])^{2}\mid x,\mathcal{D}\right] + (f(x;\mathcal{D})-\mathbb{E}[y\mid x])^{2}
\end{equation}
Here, $\mathbb{E}\left[(y-\mathbb{E}[y\mid x])^{2}\mid x,\mathcal{D}\right]$ does not depend on $\mathcal{D}$, but simply the statistical variance of $y$ given $x$. The term $(f(x;\mathcal{D})-\mathbb{E}[y\mid x])^{2}$ is considered a natural measure of effectiveness on $\mathbb{R}^{n}$ as a singular predictor of $y$. Now, for $\mathbb{E}_{\mathcal{D}}\left[(f(x;\mathcal{D})-\mathbb{E}[y\mid x])^{2}\right]$ which depends on the training set $\mathcal{D}$ in its computation, is decomposed into the form of \textit{bias-variance decomposition} terms, by derivation: 
\begin{equation}
    \begin{split}
        \mathbb{E}_{\mathcal{D}} \left[(f(x;\mathcal{D})-\mathbb{E}[y\mid x])^{2}\right] & = \underbrace{\left\{ \mathbb{E}_{\mathcal{D}}[f(x;\mathcal{D})] - \mathbb{E}[y\mid x] \right\}^{2}}_{\text{bias term}} + \underbrace{\mathbb{E}_{\mathcal{D}} \left\{(f(x;\mathcal{D})- \mathbb{E}_{\mathcal{D}}[f(x;\mathcal{D})])^{2}\right\}}_{\text{variance term}}
    \end{split}
\end{equation}

We summarize this in the following statement. 

\begin{theorem}[Bias-variance decomposition]
    Suppose the model $f(x;\mathcal{D})$ for the data $\mathcal{D}=(x_i, y_i)$ and its parameter $x$ is defined. For $y_{i}$ of the target concept's responses $y$, and consider a regression problem with the loss measure $\mathcal{M}(f,y)$ of mean squared risk, the following statement is true: \begin{equation}
        \mathbb{E}\big[\mathcal{M}(f,y)\big] = \mathcal{B}(f,y) + \mathcal{V}(f,y) + \mathbb{E}\Big[\mathbb{E} \left[((y-f(x;\mathcal{D})))^{2}\mid x, \mathcal{D}\right]\Big]
    \end{equation}
    for $\mathbb{E}[\:\cdot\mid x, \mathcal{D}]$ any expression with dependencies on $x$ and $\mathcal{D}$. The bias and variance term is subsequently expressed by 
    \begin{align}
        \mathcal{B}(f,y) = \underbrace{\left\{ \mathbb{E}_{\mathcal{D}}[f(x;\mathcal{D})] - \mathbb{E}[y\mid x] \right\}^{2}}_{\text{bias }}, \quad \mathcal{V}(f,y) =\underbrace{\mathbb{E}_{\mathcal{D}} \left\{(f(x;\mathcal{D})- \mathbb{E}_{\mathcal{D}}[f(x;\mathcal{D})])^{2}\right\}}_{\text{variance}}
    \end{align}
\end{theorem}

The above decomposition principle is often expressed into a form where there exists the intrinsic noise \cite{brown2024biasvariance}: 
\begin{equation}
    \begin{split}
        \mathbb{E}_D \left[ \mathbb{E}_{xy} \left( y - \hat{f}(x) \right)^2 \right]
        &= 
         \mathbb{E}_x \left[ \left( y^* - \mathbb{E}_D[\hat{f}(x)] \right)^2 \right]
        + \mathbb{E}_x \left[ \mathbb{E}_D \left( \hat{f}(x) - \mathbb{E}_D[\hat{f}(x)] \right)^2 \right] \\
        &+ \mathbb{E}_{xy} \left[ \left( y - y^* \right)^2 \right]
    \end{split}
    \end{equation}
For simplicity of notation, we adopt the similar form in the standard case of \cite{adlam2020understandingdoubledescentrequires}: 
\begin{equation}
    \E \qa{\hat{y}(\bfx)-y(\bfx)}^2 = \pa{\E\hat{y}(\bfx) - \E y(\bfx)}^2 + \V\qa{\hat{y}(\bfx)} + \V\q{y(\bfx)}
\end{equation}
A main common theme of criticism toward bias-variance tradeoff is the fact that the decomposition is much more general, and intrinsic for the class of \textit{mean squared loss}. However, when considering the naturalness of an error measure and then its direct applicant, mean square error on real space of sufficient support and measure comes of as a very natural choice of an error measure, at least in consideration of the regression setting; for that, we then can also define classification as another top-layer above a regression's similar real continuous space, i.e. a decision layer. Furthermore, it can also be shown \cite{brown2024biasvariance,PfauBregmanDivergence} that it also holds for the class of Bregman divergence measure. 

What does the decomposition say about model selection principle in general? In general, bias-variance is typically presented of its decomposition only, and the asymptotic prediction of its behaviours. In fact, one of the reason that it became the rule-of-thumb for ML practitioner, as well as generally statistical learning (\cite{lafon_understanding_2024} provides a quite rigorous treatment of bias-variance tradeoff in the section on statistical learning theory) solidify the trade-off as a particular model selection principle. Generally, this tradeoff can be summarized as followed: 
\begin{theorem}[Bias-variance tradeoff]
    For the expected loss of any given hypothesis $h$, the bias $\mathcal{B}(f,y)$ and variance $\mathcal{V}(f,y)$ is inversely proportional, that is, $\mathcal{B}(f,y)\propto \lambda^{-1} \mathcal{V}(f,y)$ for some proportionality $\lambda$ that may or may not be constant. In the most general case possible, $\lambda = -1$ on the entire error range. 
\end{theorem}

The tradeoff is then of inverse proportionality. Indeed, statistically, we have such tradeoff on a statistical framework in a more concrete sense. For the bias to increase, variance will increase, of which the criterion is inverse - we would like to have more bias but lower variance, and vice versa, according to such theory. 

\subsubsection{Expansion and variations}

There are problems regarding such stance. The main problem is that generally, the bias-variance measure does not totally match the overall decomposition structure. Even with irreducible errors, there are factors which make it to incorporate other error factors in for the decomposition, of which is inexplainable. There are also works, for example, \cite{domingos_unifeid_2000}, of which decompose the error to 
\begin{equation*}
        \begin{split}
            \mathrm{Error} & = \lambda_{1} \mathrm{Bias}(f,y) + \lambda_{2}\mathrm{Var}(f,y)+ \lambda_{3}\epsilon(\mathcal{D})\\ 
            & = \lambda_{1}\underbrace{\left\{ \mathbb{E}_{\mathcal{D}}[f(x;\mathcal{D})] , \mathbb{E}[y\mid x] \right\}}_{\text{bias term}} +\lambda_{2} \underbrace{\mathbb{E}_{\mathcal{D}} \left\{(f(x;\mathcal{D}), \mathbb{E}_{\mathcal{D}}[f(x;\mathcal{D})])\right\}}_{\text{variance term}} +\underbrace{\lambda_{3}\epsilon}_{\text{irreducible error}}
        \end{split}
\end{equation*}
of which the bias-variance-irreducible error are `fit' into the total error by three coefficients $\lambda_{1},\lambda_{2},\lambda_{3}$ instead. This does not only scale up the B-V-IE triplet, but also leaves out the inexplainable gap between the scaling, and the actual normal measure. 

The bias-variance decomposition and tradeoff is not general. Indeed, works, by \cite{6797087,sharma_bias-variance_2014,domingos_unifeid_2000,adlam2020understandingdoubledescentrequires,yang_rethinking_2020} and more all attempted to reconstruct and reformulate bias-variance, and it did leave out a picture of uncertainty regarding the concept. Specifically, there are certainly different way to do the bias-variance decomposition, as classified in \cite{adlam2020understandingdoubledescentrequires} paper. Some of this are particularly specialized in \textit{neural network} training, basing on the main structure of modern machine learning in neural network formalism. 
\begin{enumerate}[leftmargin=1cm,itemsep=2pt,topsep=1pt]
    \item \textbf{Semiclassical approach}: The semiclassical approach is fairly simple. We can rewrite to another variation of the classical approach to be \begin{equation}
        \underset{h\in\mathcal{H}}{\mathbb{E}[\ell(h)]} = \mathbb{E}_{x}\mathbb{E}_{\epsilon}\left[\hat{y}(x)-y(x)\right]^{2} = \mathbb{E}_{x}\left(\mathbb{E}_{\epsilon}[\hat{y}(x)]-y(x)\right)^{2} + \mathbb{E}_{x}\mathbb{V}_{\epsilon}\left[\hat{y}(x)\right]  
    \end{equation}
    and instead, consider the \textit{structural conformation} of the decomposition, or rather, taking the expectation within $\theta\in \Theta$ of the structural description of $h\in\mathcal{H}$. Then, we shall have 
    \begin{equation}
        \underset{h\in\mathcal{H}}{\mathbb{E}_{\mathrm{SC}}[\ell(h)]} = \mathbb{E}_{x}\mathbb{E}_{\theta}\mathbb{E}_{\epsilon}\left[\hat{y}(x)-y(x)\right]^{2} = \mathbb{E}_{x}\mathbb{E}_{\theta}\left(\mathbb{E}_{\epsilon}[\hat{y}(x)]-y(x)\right)^{2} + \mathbb{E}_{x}\mathbb{E}_{\theta}\mathbb{V}_{\epsilon}\left[\hat{y}(x)\right]
    \end{equation}
    Do note that the order matters, since the expectation can be constant of a given ordering. And expected value over constant is just the constant. 
    \item \textbf{Symmetric decomposition}: This approach is fairly difficult, we refer to the paper itself of its own justification on the decomposition. The symmetric decomposition relies on the shape of the input, $\mathcal{X}=\{P,D,\epsilon\}$, of which you can remove the $\epsilon$ term and retain some decomposition on $P$ and $D$ per distribution. 
\end{enumerate}
\section{Double descent}
